\section{Introduction}\label{sec:introduction}
\IEEEPARstart{F}{ace} verification compares an image of a person with a reference associated with a claimed identity. After detection and alignment, an encoder maps the face to a feature vector on which the matching decision is based. Learned representations have substantially advanced recognition~\cite{deepface,facenet,arcface}, but their identity content also creates a security concern: a disclosed vector can reveal information beyond the comparison for which it was collected. Face reconstruction from embeddings has been demonstrated in prior work~\cite{mai2019reconstruction,vec2face,shahreza2024adapter}. Replacing an exposed reference is therefore more difficult than resetting a password, because the person's biometric characteristic persists.

Cancellable biometrics introduce a replaceable transformation between the biometric features and the stored reference~\cite{ratha2001enhancing,ratha2007cancelable}. Changing its key allows a system to issue a new template after compromise. Renewal, however, is only one protection requirement. The transformed representation must also retain useful recognition performance, limit associations across records, and resist recovery of identity-bearing information. An obsolete template may be rejected while still revealing information useful for a subsequent attack. We use the terminology of biometric protection standards to distinguish these properties, without asserting formal compliance~\cite{iso24745,iso30136}.

Storage constraints make binary references attractive for edge applications such as public kiosks and smart gadgets. MobileFaceNet supplies a compact baseline encoder~\cite{chen2018mobilefacenet}; its 512-dimensional float32 output requires 2048 bytes, whereas packed 128-bit and 512-bit templates require 16 and 64 bytes. This difference becomes relevant when a device retains many references. Whether the reduction is useful nevertheless depends on the recognition performance retained and the consequences of compromise. The payload also represents only part of the resource requirement: model weights, working memory, projection operators, keys, and metadata incur additional costs.

Successive renewals can leave an adversary with several observations of the same underlying feature vector. If historical transformation operators are also exposed, hiding the current operator does not remove those observations. One-bit recovery theory explains how signs retain directional information~\cite{plan2013,jacques2013}, and pre-image and record-multiplicity attacks provide established examples of this concern~\cite{wang2021csa,li2014multiplicity}. The question examined here is how historical exposure affects renewed verification when template length, recognition utility, and encoder choice are considered together.

The study evaluates random orthogonal (RO), subsampled randomized Hadamard (SRHT), and PCA-guided RO projections on LFW and CFP-FP~\cite{huang2007lfw,sengupta2016cfp}. It compares MobileFaceNet with a ResNet50 ArcFace checkpoint under matched inputs within each encoder experiment. Official benchmark recognition and custom attack protocols are reported separately. The screened CFP attack pool contains 5704 crops from 500 identities. LFW attack results use MobileFaceNet; the two-encoder attack comparison is limited to CFP.


Renewal changes the protected reference and its operator while leaving the encoder's identity information in place. A version gate can invalidate the old record at the system interface, yet an attacker may continue to use its measurements offline. Direct replay consequently tests only whether the old bits still match; it does not test whether those bits help construct a new submission. Our evaluation treats these as separate questions and assigns a distinct outcome measure to each.

The evaluation also distinguishes the information available to each adversary. A template-only attacker attempts to associate binary records without knowing their operators. An inversion attacker receives the compromised operators and can exploit their signed constraints on the feature vector. Weak discrimination in the first setting cannot establish resistance in the second. Keeping the fresh operator unavailable during reconstruction tests whether the information retained from earlier transformations is sufficient to support acceptance after renewal.

Encoder choice complicates the interpretation of attack acceptance. A checkpoint that better matches genuine images may also match a reconstructed embedding more readily, even when reconstruction cosine changes little. Acceptance thus depends on both the recovered information and the recognition pipeline through which it is assessed. We align images, key schedules, and evaluation rules across the two encoders, while fitting preprocessing and thresholds separately on the permitted data. This compares complete frozen pipelines; it neither isolates network architecture from training nor assumes equal device cost.

The statistical analysis respects the separation between benchmark recognition and learned attacks. Official folds provide recognition estimates, whereas custom partitions keep fitting identities apart from attack-test identities. Within the attack evaluation, the same people recur across keys, exposure counts, image directions, and overlapping partitions. Resampling identities with shared weights retains these dependencies in the uncertainty estimates and paired contrasts. The intervals describe identity variation conditional on the fitted pipelines, rather than uncertainty over new training runs or a guarantee for other populations.

The first empirical contribution is a multiple-exposure security analysis. The evaluation separates unchanged-template replay, template-only linkage, and reconstruction with known old operators at $K=1,2,5,10$ exposures. The original LFW study gives \LfwBaselineAsr{} calibrated cross-image fresh-key attack success at ten exposures. In the matched CFP comparison, the corresponding ranges at 128 bits are \MbfSmallAsr{} for MobileFaceNet and \RnetSmallAsr{} for ResNet50. At 512 bits they are \MbfLargeAsr{} and \RnetLargeAsr{}, respectively. The new operator remains hidden from reconstruction. Identity-bootstrap intervals and paired comparisons support the interpretation within the tested protocols.
The second empirical contribution connects storage, recognition, and security for intended edge use. The evaluation relates 16-byte and 64-byte references to measured recognition and attack outcomes. These payloads are 128 and 32 times smaller than the raw embedding. In the matched CFP comparison, ResNet50 protected TAR is \RnetSmallTar{} at 128 bits and \RnetLargeTar{} at 512 bits, but reconstruction remains effective. The contribution is evidence for choosing compact references in intended edge settings; on-device latency, memory, and energy are not measured.

The contribution is the combined empirical assessment of exposure count, template length, and encoder choice. Neither the projection families nor the possibility of multiple-record attacks is claimed as new. Throughout the inversion experiments, the submitted object is an estimated embedding at a digital verification interface. The experiments do not establish face-image synthesis or camera-level impersonation.


Section~\ref{sec:related} reviews related work. Section~\ref{sec:methodology} defines the protection methods, attacker access, and storage accounting. Section~\ref{sec:setup} describes the datasets, encoder comparison, partitions, and statistical evaluation. Section~\ref{sec:results} presents recognition, renewal, linkage, reconstruction, and two-encoder results. Section~\ref{sec:discussion} discusses the security findings, edge implications, and limitations. Section~\ref{sec:conclusion} concludes the paper and identifies future work.
